\documentclass{article}

\usepackage[sglblindworkshop]{template/neurips_2026}
\workshoptitle{AI for Peace Workshop}
\makeatletter
\renewcommand{\@noticestring}{Prepared for the NeurIPS 2026 AI for Peace Workshop.}
\makeatother

\usepackage[utf8]{inputenc}
\usepackage[T1]{fontenc}
\usepackage{hyperref}
\usepackage{url}
\usepackage{booktabs}
\usepackage{amsmath}
\usepackage{graphicx}
\usepackage{microtype}
\usepackage{array}
\hypersetup{colorlinks=true,citecolor=blue,urlcolor=blue}

\title{Matched Source Controls Expose a Capability Boundary in Appeal Review}
\author{Rafael Gardos\\\texttt{gardos.rafael@gmail.com}}

\begin{document}
\nolinenumbers

\maketitle

\begin{abstract}
Plausible evidence can still be inadmissible under a service policy, and apparent accuracy depends on the disposition convention used to score it. We pair 48 synthetic humanitarian appeals from unlisted sources with 24 controls from accepted sources. Under the printed policy, four local models return the exact invalid-source disposition in 0 of 384 independent reviews. An alternative \texttt{NEED\_INFORMATION} target changes the model ranking. Mistral's independent balanced exact score rises from 38.5\% to 85.4\%, and Phi's rises from 25.0\% to 62.5\%. Qwen and Gemma remain constant acceptors. GPT-5.6 Sol separates all 96 unlisted records from all 48 accepted controls under both review contexts. It returns the printed disposition with 100\% balanced exact accuracy, while the alternative convention scores the same outputs at 50\%. Every frontier answer is strictly parsed and fully grounded. Earlier-rationale exposure raises local Mistral's false eligibility by 35.42 points, with a 12-base-request interval of [29.17, 41.67]. The cases support no operational aid decision. We conclude with a three-step Matched Admissibility Audit for model-supported review.
\end{abstract}

\section{Introduction}

An appeal channel can exist while offering little chance of correction. A reviewer may accept a current-looking record even when the governing policy does not recognize its source. This concern matters when a language model helps staff sort requests or draft replies. It also applies when a model prepares a record for later human review. Human presence at the end of a workflow does not reveal whether the model has already narrowed the options that reach that person.

Humanitarian accountability gives this problem a direct connection to peace work. The Core Humanitarian Standard asks organizations to support crisis-affected people in ways that respect rights and dignity. It also calls for people and communities to hold those organizations to account \citep{chs2024standard}. Research on artificial intelligence in humanitarian action warns that technical benefits arrive with risks to rights and vulnerable populations \citep{beduschi2022humanitarian}. An accurate model response cannot create an accessible complaint channel or supply accountable authority. It can still affect the information and recommendation that enter such a channel.

We study one bounded behavior in that process. The primary experiment asks whether a model discriminates between accepted and unlisted sources. Exact disposition remains co-primary because it determines the procedure that follows. We also ask whether a case-specific denial rationale changes source discrimination. The rationale has no evidentiary status under the supplied rules.

The primary set contains 48 semantic invalid-source cases across four of the eight fictional humanitarian service families. Four plausible source types replace an accepted source while the rest of the record remains current. The matched control uses the same 12 base requests and instantiates both source types that the policy accepts. Each case appears in two meaning-preserving orders. A source audit found no appeal item or saved model review reused elsewhere. Similar service families were treated as topic overlap. We test independent review and prior-rationale exposure with four pinned local instruction models. An earlier 96-case set supplies a secondary audit of valid, invalid, incomplete and conflicting evidence.

The experimental unit is one appeal review under a named source rule. No field-verifier judgment, memory decision or readability rewrite from the other studies enters an AppealShift score.

The central result is heterogeneous. Qwen and Gemma are constant acceptors across the matched source classes. Phi rejects every unlisted source and accepts only half of the valid controls. Mistral discriminates best, yet asks for information where C2 requires rejection. None of 384 independent reviews returns the exact invalid-source disposition. The earlier rationale makes Mistral less selective. The matched design shows why a zero or a rejection rate alone cannot identify source-rule capability.

The contribution combines a controlled admissibility audit with a procedural requirement for humanitarian accountability. AppealShift holds the request and eligibility fact fixed while crossing accepted and unlisted source types. It then varies prompt order and adds or removes denial context. An organization considering AI support for appeals should test both sides of a source rule. It should inspect the disposition chosen for rejected evidence instead of counting every non-acceptance as correct.

Procedural completion can conceal weak evidence handling. AppealShift therefore reports the disposition apart from evidence grounding across matched review contexts.

\section{Related work and claim boundary}

\subsection{Contestability and humanitarian accountability}

Contestability concerns how a person can challenge a decision and how an institution can correct error. The concept has several perspectives and cannot be reduced to an explanation interface \citep{lyons2021contestability}. Recent work separates contestability from recourse and identifies evidence that can make a decision indefensible \citep{freiesleben2026contestability}. AppealShift does not propose a new definition or evidence taxonomy. It tests one model behavior that a contestable process may need to detect.

The distinction matters in humanitarian settings. Accountability depends on organizational practice and community influence, not model accuracy alone \citep{chs2024standard}. The cases in this study concern intake and scheduling. Other cases cover referrals and callbacks. They avoid allocation of food or shelter. Medical treatment and legal status also remain outside the benchmark. This boundary keeps the experiment tied to humanitarian service review without turning a synthetic policy into a claim about a real institution.

\subsection{Appeals, admissibility and evidence updating}

Legal datasets already connect initial and appellate decisions. AppealCase contains 10,000 matched civil-case pairs with reversal and new-information annotations \citep{huang2025appealcase}. AppellateGen studies judgment generation over 7,351 case pairs and evidentiary updates \citep{yang2026appellategen}. These datasets are larger and closer to legal practice. They study what changes when new evidence arrives. AppealShift asks whether a reviewer distinguishes evidence content from the source rule that makes the evidence admissible.

Production review systems supply another close comparison. Reviewing the Reviewer represents evidence and verification actions in e-commerce appeals, then records the resulting decision \citep{du2026reviewer}. Its workflow can request missing information and retrieve prior correction patterns. AppealShift makes no novelty claim for checklists or information requests. Its matched control changes only whether the stated record type appears on the policy whitelist. The rationale comparison then tests whether prior case text changes that source discrimination.

Recent agent-memory evaluations show that stored content can acquire authority after source boundaries are compressed, and that provenance can be laundered through persistent memory \citep{zhan2026authority,xu2026laundering}. Those studies concern consolidation and later retrieval. AppealShift tests a current review prompt in which the source rule is explicit and the evidentiary fact is matched. The overlap is the need to keep content separate from the authority of its source.

\subsection{Prior text can change later judgment}

Rationales can alter judgments without adding evidence. Pro and con rationales shift human and model plausibility ratings on commonsense questions \citep{palta2026arguments}. Model beliefs can also distort ratings of persuasive quality \citep{zahraei2026beliefs}. Numerical anchoring changes quantitative judgments in ways that depend on model confidence and post-training \citep{owusu2026anchoring}.

A recent verifier study is especially close in design. A completed audit and repair episode about another item shifted false-alarm thresholds on a byte-identical current task \citep{mazaheri2026audit}. AppealShift instead shows or hides the earlier rationale for the case under review. The current evidence stays fixed, and the output is one of four categorical service dispositions. This distinction supports a narrow gap. Prior work establishes context sensitivity in other judgment tasks. It does not answer how case-specific denial reasoning affects these matched fictional humanitarian reviews.

Authority and dialogue framing provide two further boundaries. Sycophants in the Courtroom finds that models over-trust authoritative but false legal information and react to superficial formatting cues \citep{molfetta2026sycophants}. Rabbani et al. show that minimal conversational framing and a rebuttal can alter model judgment \citep{rabbani2026framing}. These studies establish susceptibility to authority and conversational pressure. AppealShift makes no general claim for either phenomenon. Its manipulated text is the earlier denial rationale for the same case, which has no evidentiary status under the supplied policy. The source comparison instead holds the eligibility fact fixed and changes whether the named record type is accepted.

\section{AppealShift}

\subsection{Scope and cases}

All records are synthetic. No real person or private record appears. The agencies and places are fictional. So are the policies. The eight service families include family-contact appointments and accessible transport scheduling. Replacement-document intake and interpretation booking form another pair. Remote check-ins and legal-information referrals appear as well. Shelter-maintenance referrals and complaint callbacks complete the coverage. These workflows test whether a fixed procedural relation survives vocabulary changes. They do not sample actual humanitarian institutions.

Each case has four policy clauses. Clause C1 states the eligibility fact. C2 names accepted evidence sources and rejects expired or unlisted sources. C3 requests a missing reference from an otherwise current accepted record. C4 sends conflicting accepted records to a human reviewer. The fixed decision order applies C4 first and C3 next. C2 follows, with C1 last.

Four evidence states map to exact dispositions. Valid evidence supports \texttt{ELIGIBLE}. Invalid evidence fails C2 and targets \texttt{INELIGIBLE}. An incomplete record targets \texttt{NEED\_INFORMATION}. Conflicting accepted records target \texttt{HUMAN\_REVIEW}. There are 24 semantic cases per state. A policy-first form and a record-first form yield 192 evaluation rows. A disjoint 16-row development set is excluded from every reported estimate.

This label map belongs to the fictional policy. The humane service response remains an open design choice. A practitioner could reasonably specify \texttt{NEED\_INFORMATION} for an unlisted but current-looking source so the person can provide an accepted record. That alternative would change C2 and its target. AppealShift measures adherence to the rule printed in each case. The benchmark makes no general claim that \texttt{INELIGIBLE} is the best policy.

The shared decision-order paragraph was added after the first development run exposed ambiguous precedence in the natural-language clauses. No evaluation response existed at that point. Version 1 development files remain in the archive. The data and labels were retained. The conditions and confirmatory contrast were retained as well.

\subsection{Review conditions}

Independent review hides the earlier disposition and rationale. The model receives the fictional policy with the current record. Appeal evidence appears in the same prompt. Prior-rationale review adds the earlier \texttt{INELIGIBLE} disposition with its case-specific explanation. It adds no new evidence.

The evidence-checklist workflow includes the denial context and asks for a clause-level evidence check before the structured answer. Commit-first review begins with an independent disposition. A second model call receives that preliminary answer and the earlier denial rationale, then returns the final review. Only final answers enter the main endpoint. Preliminary answers support an update analysis.

Every answer must be one JSON object with a disposition and policy clause. It also carries the evidence identifiers and a short reply. The main scorer accepts only the four specified dispositions. It verifies exact keys and reports strict parsing separately from recovered parsing. A second scorer was written without importing the primary scoring code.

\subsection{Recorded plausible-source experiment}

The extension contains 48 semantic invalid-evidence cases. Twelve base cases cross four plausible but unlisted source types. Two are a portal export and a forwarded screenshot. The others are a staff chat and a coordination digest. The fictional policy lists other accepted sources and states that an unlisted source is invalid even when its contents look current.

Each semantic case has policy-first and record-first forms, giving 96 surface records. Independent review and prior-rationale review create 192 decisions per model and 768 overall. Every cell contains 96 decisions. The target is always \texttt{INELIGIBLE} under clause C2. Source identity and surface order are balanced within each model-condition cell. The generator was refactored to expose a pure build function before rows were written. A separate validator checked the cross-product and labels. It checked source exclusion and phrase removal in separate passes. The dataset digest was checked as well.

\subsection{Matched valid-source controls}

The control uses the same 12 base requests. C2 names two accepted source types for each request family, and the generator creates one record from each type. It copies the request identifier and C1 fact from the matched invalid case. The fixed evidence sentence states that the current accepted record identifies the request and directly confirms C1. The target is \texttt{ELIGIBLE} under C1 with evidence E1.

The 24 semantic controls appear in both surface orders. Two review conditions and four models produce 384 decisions. The released artifact includes the control protocol and complete data. The control construction does not alter the invalid-source records.

\subsection{Models and execution}

The model set contains Qwen3 4B Instruct 2507 and Phi-4-mini Instruct. Gemma 3 4B Instruct and Mistral 7B Instruct v0.3 complete the grid. Each model revision is pinned by repository commit. All runs use local four-bit MLX weights with greedy decoding. The output limit is 180 new tokens. Native chat templates are retained. The Mistral adapter folds the system text into the first user message because its template rejects a system role.

Each model produces 768 final reviews on the earlier four-state grid. That evaluation contains 3,072 final reviews and 768 preliminary commitment calls. Every model, condition, evidence-state cell contains 48 records. The primary plausible-source set adds 192 reviews per model, and the matched control adds 96. Raw outputs are retained without result-dependent removal.

Two earlier checks remain exploratory. The first reruns all 48 valid-case surface variants under independent and prior-rationale review with matching BF16 Qwen and Gemma artifacts. This adds 192 reviews. The second is the 12-case plausible-source slice that motivated the prospective expansion. It adds 192 reviews across the original four-bit models and is kept separate from the new primary set.

After the local source comparison was sealed, GPT-5.6 Sol reviewed all 96 unlisted records and 48 matched controls under independent and prior-rationale context. The 288 calls use seed 20260827 with reasoning disabled and provider-managed temperature. The runner records the requested and returned model slugs, provider, token use and generation identifier. This post-audit replication tests whether the local source failures remain at a stronger endpoint. It does not estimate frontier-model behavior in general.

\section{Analysis}

Balanced source discrimination and exact invalid-source disposition are co-primary descriptive endpoints for the plausible-source experiment. Balanced discrimination averages valid-source sensitivity and unlisted-source specificity. Specificity treats any non-\texttt{ELIGIBLE} answer as rejection, while exact accuracy still requires the policy's stated disposition. False eligibility and fully grounded correctness are co-reported diagnostics. Fully grounded correctness also requires the cited clause and decisive evidence identifiers to agree with the deterministic target.

Mistral supplies the only false-eligibility change. Its interval uses 20,000 bootstrap draws clustered on the 12 base requests. Each sampled base carries both surface forms and all four unlisted source variants under both conditions. The seed is 20260825. With only 12 clusters, percentile coverage is approximate. The interval conditions on this tested artifact and does not estimate a model population.

The sole confirmatory contrast in the earlier four-state experiment remains prior-rationale minus independent-review accuracy on valid appeals. It is retained as a secondary fixed-artifact result rather than promoted above the source-validity test.

The earlier valid-appeal contrast pools four fixed models and both surface forms. Only Mistral moves, so Table~\ref{tabvalid} reports its model-specific interval rather than a pooled interval.

Every model-specific result is exploratory. The two additional workflows are exploratory as well. Results on the other evidence states receive the same treatment. Their intervals are unadjusted, with no multiplicity correction.

Validation reconstructs every prompt and score, then checks model revisions and dataset hashes. The primary and independent scorers agree on every local final record. Eight of 3,072 final outputs in the earlier grid were unparseable. All 1,152 local source-comparison reviews and all 288 frontier reviews parsed.

\section{Results}

\subsection{Matched sources expose three different failures}

Independent-review accuracy is zero across 384 invalid-source decisions. Fully grounded accuracy is also zero. Table~\ref{tabconfusion} shows what that pooled zero conceals. Qwen and Gemma choose \texttt{ELIGIBLE} every time. Mistral asks for information on 90 of 96 cases. Phi sends 72 to information requests and 24 to human review. The latter answers reject the unlisted record but do not apply C2's required \texttt{INELIGIBLE} disposition.

\begin{table}[t]
\caption{Independent invalid-source reviews by predicted disposition. Each row has 96 decisions. E, I, N and H denote \texttt{ELIGIBLE}, \texttt{INELIGIBLE}, \texttt{NEED\_INFORMATION} and \texttt{HUMAN\_REVIEW}. $\dagger$ marks a constant acceptor.}
\label{tabconfusion}
\centering
\small
\begin{tabular}{lrrrr}
\toprule
Model & E & I & N & H \\
\midrule
Qwen3 4B$^\dagger$ & 96 & 0 & 0 & 0 \\
Mistral 7B & 6 & 0 & 90 & 0 \\
Phi-4-mini & 0 & 0 & 72 & 24 \\
Gemma 3 4B$^\dagger$ & 96 & 0 & 0 & 0 \\
\bottomrule
\end{tabular}
\end{table}

The matched controls identify whether those behaviors discriminate source class. Qwen and Gemma accept all 48 valid controls and all 96 unlisted records. Their balanced discrimination is 50\%, the constant-acceptance value. Phi rejects every unlisted source but accepts only 24 of 48 valid records. Mistral accepts 37 valid records and rejects 90 unlisted records. It reaches 85.4\% balanced discrimination while never returning the exact invalid disposition.

\begin{table}[t]
\caption{Matched source discrimination in percent. Sensitivity is acceptance of 48 valid controls. Specificity is non-acceptance of 96 unlisted records. Exact is correct invalid-source disposition accuracy. Balanced averages sensitivity and specificity. $\dagger$ marks a constant acceptor.}
\label{tabmatched}
\centering
\small
\begin{tabular}{llrrrr}
\toprule
Model & Review & Sens. & Spec. & Exact & Balanced \\
\midrule
Qwen3 4B$^\dagger$ & Independent & 100.0 & 0.0 & 0.0 & 50.0 \\
          & Rationale   & 100.0 & 0.0 & 0.0 & 50.0 \\
Mistral 7B & Independent & 77.1 & 93.8 & 0.0 & 85.4 \\
           & Rationale   & 100.0 & 58.3 & 0.0 & 79.2 \\
Phi-4-mini & Independent & 50.0 & 100.0 & 0.0 & 75.0 \\
           & Rationale   & 20.8 & 100.0 & 25.0 & 60.4 \\
Gemma 3 4B$^\dagger$ & Independent & 100.0 & 0.0 & 0.0 & 50.0 \\
           & Rationale   & 100.0 & 0.0 & 0.0 & 50.0 \\
\bottomrule
\end{tabular}
\end{table}

The exact-disposition result depends on the policy convention. Table~\ref{tabconventions} keeps the valid target at \texttt{ELIGIBLE} and changes only the invalid target used for rescoring. No convention makes every model correct. The alternative changes the model ranking because Mistral and Phi often request information, while the two constant acceptors cannot benefit.

\begin{table}[t]
\caption{Balanced exact-disposition accuracy in percent under two invalid-source targets. Each score averages a 48-record valid cell and a 96-record invalid cell. I uses \texttt{INELIGIBLE}. N uses \texttt{NEED\_INFORMATION}.}
\label{tabconventions}
\centering
\small
\begin{tabular}{lrr@{\hspace{1.3em}}rr}
\toprule
& \multicolumn{2}{c}{Independent} & \multicolumn{2}{c}{Rationale} \\
\cmidrule(lr){2-3}\cmidrule(lr){4-5}
Model & I & N & I & N \\
\midrule
Qwen3 4B$^\dagger$ & 50.0 & 50.0 & 50.0 & 50.0 \\
Mistral 7B & 38.5 & 85.4 & 50.0 & 79.2 \\
Phi-4-mini & 25.0 & 62.5 & 22.9 & 29.2 \\
Gemma 3 4B$^\dagger$ & 50.0 & 50.0 & 50.0 & 50.0 \\
\bottomrule
\end{tabular}
\end{table}

The source-type breakdown rules out one troublesome document label as the sole cause. Exact accuracy remains zero for each of the four unlisted types. Independent false eligibility ranges from 50.00\% to 56.25\% across those types.

\subsection{A stronger reviewer resolves the printed source rule}

GPT-5.6 Sol accepts all 48 matched valid controls and rejects all 96 unlisted records under both review contexts. Balanced source discrimination is therefore 100\%. Every invalid review returns \texttt{INELIGIBLE} with C2 and E1, while every valid review returns \texttt{ELIGIBLE} with C1 and E1. Strict schema and full-grounding rates are also 100\%. Showing the earlier rationale changes none of the 288 dispositions.

The dual-convention result remains necessary. Under the printed C2 target, frontier balanced exact accuracy is 100\%. Rescoring the same invalid outputs against \texttt{NEED\_INFORMATION} lowers it to 50\%, because every valid control remains correct and every invalid disposition becomes wrong. The matched task is therefore solvable at this endpoint, while the local failures mark a capability boundary on these prompts. A policy owner must still choose and publish the disposition convention before an audit score has procedural meaning.

\subsection{A denial rationale raises false eligibility}

With the earlier rationale visible, pooled disposition accuracy rises to 6.25\%. Phi supplies the entire change. Its 24 correct dispositions still cite the wrong clause or evidence, leaving fully grounded accuracy at zero. A fixed five-case record audit found substantive mismatches rather than formatting variants. Every case cites C1, C3 or C4 instead of C2. One also omits E1. The other three models remain at zero disposition accuracy.

Pooled false eligibility rises from 51.56\% to 60.42\%, a descriptive average change of 8.85 points.\footnote{The displayed condition rates are rounded separately and differ by 8.86 points. The paired estimate uses unrounded counts and equals 8.854 points, reported as 8.85.} Mistral accounts for the entire movement. It rises from 6.25\% to 41.67\%, a 35.42-point increase with a 12-base-request interval from 29.17 to 41.67. The other three models are unchanged. Under rationale exposure, false eligibility reaches 51.04\% for policy-first prompts and 69.79\% for record-first prompts. The order gap is descriptive, since wording beyond block order also changes with the surface form.

The larger invalid-source result reproduces zero independent exact accuracy and balances four source types. The matched controls prevent that zero from being read as one common failure. The original slice remains in the artifact as the observation that motivated the larger test.

\subsection{Secondary evidence-state checks}

On the earlier valid appeals, pooled local accuracy rises from 83.85\% to 90.62\% after rationale exposure. Mistral supplies all 13 improvements and rises 27.08 points [16.67, 39.58]. Omitting it reduces the pooled gain to zero. A BF16 check reproduces the Qwen and Gemma ceiling cells exactly, so those two cells cannot reveal a condition effect at either precision.

\subsection{Accuracy hid changes in error type}

The rationale changed outcomes outside the confirmatory state. Invalid-case accuracy rose from 9.38 to 24.48 percent. False eligibility on those same records also rose from 49.48 to 61.46 percent. Phi supplied most of the new correct invalid decisions. Mistral introduced the adverse shift in false eligibility, moving from zero to 50 percent. Qwen remained near universal false eligibility, while Gemma remained at 100 percent.

\begin{table}[t]
\caption{State-specific disposition accuracy in percent. The final column records the safety-relevant response for that state.}
\label{tabstates}
\centering
\small
\begin{tabular}{lrrp{0.37\linewidth}}
\toprule
State & Independent & Rationale & Additional observation \\
\midrule
Valid & 83.85 & 90.62 & Failure to correct fell from 16.15 to 9.38 \\
Invalid & 9.38 & 24.48 & False eligibility rose from 49.48 to 61.46 \\
Incomplete & 50.00 & 35.94 & Information requests fell by 14.06 points \\
Conflict & 11.46 & 13.54 & Human review rose by 2.08 points \\
\bottomrule
\end{tabular}
\end{table}

Incomplete-case accuracy fell from 50.00 to 35.94 percent. Appropriate information requests fell by the same 14.06 points. Conflict routing rose from 11.46 to 13.54 percent. Across all evidence states, rationale accuracy was 41.15 percent and independent accuracy was 38.67 percent. Fully grounded correctness moved in the other direction, from 27.34 percent to 26.17 percent. Schema validity improved from 79.30 to 87.11 percent, while evidence correctness fell from 79.04 to 73.96 percent.

Independent-review accuracy was only 9.38 percent on invalid records, and appropriate conflict routing was 11.46 percent. These floor effects make the deltas unsuitable as evidence of task capability at this model scale. They still show which wrong action appeared. The same exposure can improve the target label in one state and increase a consequential error in another.

The exploratory 12-case slice had zero independent accuracy across 96 reviews and motivated the larger experiment. It is not pooled. The evidence checklist changes overall accuracy by $-0.91$ points $[-2.73,0.91]$. Commit-first review changes 194 of 768 preliminary answers, with 56 corrections and 50 harmful changes relative to those preliminary answers. Against prior-rationale final review, it loses 58 correct decisions and raises false eligibility on invalid evidence by 10.42 points. Neither workflow supplies a reliable prompt fix. Full state, workflow and surface breakdowns remain in the artifact.

\section{Accountability requirement for humanitarian review}

The matched sources make admissibility the headline. Qwen and Gemma ignore whether C2 accepts the source. Phi routes many valid and invalid cases away from a decision. Mistral usually distinguishes source classes, yet calls an unlisted record missing information. Each behavior requires a different repair.

The exact disposition determines the procedure recommended inside this policy. A request for information seeks another record. An ineligible decision exposes the source-rule rejection, while human review invokes an accountable route. A real service may choose another mapping. Binary acceptance would hide these differences.

The earlier rationale has no evidentiary status in the fictional policy, yet changes source discrimination. For Mistral, valid-source sensitivity reaches 100\% while unlisted-source specificity falls to 58.3\%. Phi moves in the opposite direction on valid controls. This effect remains artifact-specific.

We call the resulting procedure the Matched Admissibility Audit.
\begin{enumerate}
\item Pair each unlisted-source record with accepted-source controls that preserve the request and eligibility fact.
\item Score source discrimination and the exact four-way disposition under every policy convention the service is prepared to defend.
\item Repeat the matched cells across prompt order and each deployed model revision, then retain the full disposition counts for accountable review.
\end{enumerate}

The workflow comparisons give no simple prompt fix. The checklist had no reliable overall gain, and prior commitment made performance worse. These findings apply to the exact prompts and model revisions in the study.

The source-validity concern would weaken if stronger reviewers accepted the matched valid records while rejecting plausible unlisted sources with the C2 disposition and evidence. A replication in which broader source wording removed the error would locate it in the present templates. The rationale effect would fail if repeated runs removed Mistral's false-eligibility increase.

For peace and humanitarian work, this yields a concrete procedural requirement. New evidence must remain able to change the record that reaches an accountable reviewer. Before prior case reasoning enters a model, the organization should show that it does not suppress requests for missing evidence or redirect cases into unacceptable dispositions. The test should be repeated across prompt order and each model revision in use. Passing that audit still cannot replace an accessible complaint mechanism or determine whether remedy arrives in time.

\section{Limits and ethical use}

The benchmark is in English and uses short templates. Repeated wording may cue models. The two surface orders test one wording change. Translation and longer histories remain outside scope. Four selected source types cannot represent open-ended document variation. The frontier result covers one routed model endpoint and does not represent larger review systems as a class.

The fictional workflows avoid direct allocation decisions, which narrows ecological validity. Real services have local rules. Affected people may lack accepted evidence or face danger when using a complaint channel. Those conditions do not appear here. The study covers four local models with greedy four-bit decoding. It does not represent proprietary systems or larger reasoning models.

Exact labels reproduce scores, not policy legitimacy. AppealShift is an audit fixture and must never adjudicate requests or rank people. Real use requires accountable human review and an accessible remedy channel.

No human participant or private record was used. Every case comes from a fictional template. The score-blind packet has no external labels, so no completed human annotation study is claimed.

\paragraph{Reproducibility.} Code, data, protocol records, raw outputs and audits are public at \url{https://github.com/xyz-rafael-xyz/appealshift-neurips-2026}. The public history does not establish when any protocol or analysis plan was written relative to generation. A separate program reproduces the earlier-grid claims, and a deterministic analysis verifies the 288 frontier reviews. The one-page abstract renders the locally validated form payload supplied with the artifact. Submission through the live form remains an author action. This supporting paper uses a longer technical abstract.

\clearpage
\bibliographystyle{plainnat}
\bibliography{references}

\appendix

\section{Model revisions}

\begin{table}[h]
\caption{Local model artifacts used in every reported run.}
\centering
\small
\begin{tabular}{p{0.21\linewidth}p{0.38\linewidth}p{0.29\linewidth}}
\toprule
Model & Repository & Revision \\
\midrule
Qwen3 4B & \texttt{mlx-community/}\newline\texttt{Qwen3-4B-Instruct-}\newline\texttt{2507-4bit} & \texttt{50d427756c6b1b2fe0c0}\newline\texttt{a10f67fbda1fc8e82c1b} \\
Phi-4-mini & \texttt{mlx-community/}\newline\texttt{Phi-4-mini-instruct-}\newline\texttt{mlx-4Bit} & \texttt{d848c30f6d5419b98924}\newline\texttt{33cf6b1062626d15340e} \\
Gemma 3 4B & \texttt{mlx-community/}\newline\texttt{gemma-3-text-4b-it-}\newline\texttt{4bit} & \texttt{4f665a4c50ecfe4ecdc3}\newline\texttt{4056ab52fe3e3c4abf9e} \\
Mistral 7B & \texttt{mlx-community/}\newline\texttt{Mistral-7B-Instruct-}\newline\texttt{v0.3-4bit} & \texttt{a4b8f870474b0eb527f4}\newline\texttt{66a03fbc187830d271f5} \\
\bottomrule
\end{tabular}
\end{table}

\section{Dataset balance}

Each of the eight service families contributes three semantic cases to each target disposition. This creates 12 semantic cases per family and 96 in total. Each semantic case has two surface forms. Seventy-two cases have one decisive evidence identifier. The 24 conflict cases require two.

\begin{flushleft}
The evaluation file has SHA256 \texttt{5a7a74533c614dcdd244003bf9411d88}\allowbreak\texttt{3237daac818731625d3229ea4c7ee9cf}.

The development file has SHA256 \texttt{551e1040dbe6d6fd9f365d4267115eb1}\allowbreak\texttt{24fa2a7a6a2c2740b728de5554038a24}.
\end{flushleft}

\section{Validation record}

The full-run validator checks prompt reconstruction and exact model revisions. It then checks dataset hashes and row coverage. Score agreement has a separate check. The independent result checker reads raw JSONL without importing the main analysis. It passed 1,080 of 1,080 checks. Eight final outputs were unparseable under the declared recovery parser. No record was removed.

The archive also contains a 384-record score-blind audit packet selected across 128 strata using identifiers only. It is included for later external review. Its status remains unlabelled, and no paper claim treats it as a human audit.

\end{document}
